\uuid{wtlR}
\titre{Noyau/Image/Injection/Surjection/Théorème du rang/Polynôôme}
\auteur{hanani ; mzari}
\organisation{exo7}
\licence{CC BY-NC-SA 4.0}
\datecreate{2020-01-20}
\niveau{L1}
\module{Algèbre}
\chapitre{Application linéaire}
\sousChapitre{Définition}
\difficulte{3}
\theme{Noyau et image, Noyau, Image, Injection, Surjection, Théorème du rang, Polynôôme}
\exo7qcmid{lille-2.applications-lineaires.022}

\contenu{
\begin{question}
\qclassification{107.01}
\qtags{motcle=Noyau/Image/Injection/Surjection/Théorème du rang/Polynôôme}
On considère $\Rr_3[X]$, l'espace des polynômes à coefficients réels de degré $\le 3$, muni de sa base canonique ${\cal{ B}}= \{1,X,X^2,X^3\} $ et $f$ l'endomorphisme de $\Rr_3[X]$
défini par :
$$f(1)=X,\;  f(X)=1+X,\; f(X^2)= (X-1)^2,\; f(X^3)=(X-1)^3.$$
Quelles sont les assertions vraies ?
\begin{answers}
\bad{$\dim \ker f =1$}.
\good{$f$ est injective}.
\bad{$f$ n'est pas injective}.
\good{$\mbox{rg} (f)=4$}.
\end{answers}
\begin{explanations}  Soit $P= aX^3+bX^2+cX+d,\, a,b,c,d \in \Rr $. On a :
$$P\in \ker f\Leftrightarrow a(X-1)^3+b(X-1)^2+(c+d)(X-1)+2c+d=0.$$ Comme $\{1,X-1,(X-1)^2,(X-1)^3\}$ est une famille libre, on déduit que $P=0$. Donc $\dim \ker f=0$ et $f$ est injective.
\vskip0mm
D'après le théorème du rang, $\mbox{rg}(f)=\dim \Im f=4$ et comme $\Im f$ est un sous-espace vectoriel de $\Rr_3[X]$, $\Im f=\Rr_3[X]$ et donc $f$ est surjective.
\end{explanations}
\end{question}
}
